% Options for packages loaded elsewhere
\PassOptionsToPackage{unicode}{hyperref}
\PassOptionsToPackage{hyphens}{url}
\PassOptionsToPackage{dvipsnames,svgnames,x11names}{xcolor}
%
\documentclass[
  letterpaper,
  DIV=11,
  numbers=noendperiod]{scrreprt}

\usepackage{amsmath,amssymb}
\usepackage{iftex}
\ifPDFTeX
  \usepackage[T1]{fontenc}
  \usepackage[utf8]{inputenc}
  \usepackage{textcomp} % provide euro and other symbols
\else % if luatex or xetex
  \usepackage{unicode-math}
  \defaultfontfeatures{Scale=MatchLowercase}
  \defaultfontfeatures[\rmfamily]{Ligatures=TeX,Scale=1}
\fi
\usepackage{lmodern}
\ifPDFTeX\else  
    % xetex/luatex font selection
\fi
% Use upquote if available, for straight quotes in verbatim environments
\IfFileExists{upquote.sty}{\usepackage{upquote}}{}
\IfFileExists{microtype.sty}{% use microtype if available
  \usepackage[]{microtype}
  \UseMicrotypeSet[protrusion]{basicmath} % disable protrusion for tt fonts
}{}
\makeatletter
\@ifundefined{KOMAClassName}{% if non-KOMA class
  \IfFileExists{parskip.sty}{%
    \usepackage{parskip}
  }{% else
    \setlength{\parindent}{0pt}
    \setlength{\parskip}{6pt plus 2pt minus 1pt}}
}{% if KOMA class
  \KOMAoptions{parskip=half}}
\makeatother
\usepackage{xcolor}
\usepackage[a4paper,total=\{170mm, 257mm\},left=20mm,top=20mm]{geometry}
\setlength{\emergencystretch}{3em} % prevent overfull lines
\setcounter{secnumdepth}{-\maxdimen} % remove section numbering
% Make \paragraph and \subparagraph free-standing
\ifx\paragraph\undefined\else
  \let\oldparagraph\paragraph
  \renewcommand{\paragraph}[1]{\oldparagraph{#1}\mbox{}}
\fi
\ifx\subparagraph\undefined\else
  \let\oldsubparagraph\subparagraph
  \renewcommand{\subparagraph}[1]{\oldsubparagraph{#1}\mbox{}}
\fi

\usepackage{color}
\usepackage{fancyvrb}
\newcommand{\VerbBar}{|}
\newcommand{\VERB}{\Verb[commandchars=\\\{\}]}
\DefineVerbatimEnvironment{Highlighting}{Verbatim}{commandchars=\\\{\}}
% Add ',fontsize=\small' for more characters per line
\usepackage{framed}
\definecolor{shadecolor}{RGB}{241,243,245}
\newenvironment{Shaded}{\begin{snugshade}}{\end{snugshade}}
\newcommand{\AlertTok}[1]{\textcolor[rgb]{0.68,0.00,0.00}{#1}}
\newcommand{\AnnotationTok}[1]{\textcolor[rgb]{0.37,0.37,0.37}{#1}}
\newcommand{\AttributeTok}[1]{\textcolor[rgb]{0.40,0.45,0.13}{#1}}
\newcommand{\BaseNTok}[1]{\textcolor[rgb]{0.68,0.00,0.00}{#1}}
\newcommand{\BuiltInTok}[1]{\textcolor[rgb]{0.00,0.23,0.31}{#1}}
\newcommand{\CharTok}[1]{\textcolor[rgb]{0.13,0.47,0.30}{#1}}
\newcommand{\CommentTok}[1]{\textcolor[rgb]{0.37,0.37,0.37}{#1}}
\newcommand{\CommentVarTok}[1]{\textcolor[rgb]{0.37,0.37,0.37}{\textit{#1}}}
\newcommand{\ConstantTok}[1]{\textcolor[rgb]{0.56,0.35,0.01}{#1}}
\newcommand{\ControlFlowTok}[1]{\textcolor[rgb]{0.00,0.23,0.31}{#1}}
\newcommand{\DataTypeTok}[1]{\textcolor[rgb]{0.68,0.00,0.00}{#1}}
\newcommand{\DecValTok}[1]{\textcolor[rgb]{0.68,0.00,0.00}{#1}}
\newcommand{\DocumentationTok}[1]{\textcolor[rgb]{0.37,0.37,0.37}{\textit{#1}}}
\newcommand{\ErrorTok}[1]{\textcolor[rgb]{0.68,0.00,0.00}{#1}}
\newcommand{\ExtensionTok}[1]{\textcolor[rgb]{0.00,0.23,0.31}{#1}}
\newcommand{\FloatTok}[1]{\textcolor[rgb]{0.68,0.00,0.00}{#1}}
\newcommand{\FunctionTok}[1]{\textcolor[rgb]{0.28,0.35,0.67}{#1}}
\newcommand{\ImportTok}[1]{\textcolor[rgb]{0.00,0.46,0.62}{#1}}
\newcommand{\InformationTok}[1]{\textcolor[rgb]{0.37,0.37,0.37}{#1}}
\newcommand{\KeywordTok}[1]{\textcolor[rgb]{0.00,0.23,0.31}{#1}}
\newcommand{\NormalTok}[1]{\textcolor[rgb]{0.00,0.23,0.31}{#1}}
\newcommand{\OperatorTok}[1]{\textcolor[rgb]{0.37,0.37,0.37}{#1}}
\newcommand{\OtherTok}[1]{\textcolor[rgb]{0.00,0.23,0.31}{#1}}
\newcommand{\PreprocessorTok}[1]{\textcolor[rgb]{0.68,0.00,0.00}{#1}}
\newcommand{\RegionMarkerTok}[1]{\textcolor[rgb]{0.00,0.23,0.31}{#1}}
\newcommand{\SpecialCharTok}[1]{\textcolor[rgb]{0.37,0.37,0.37}{#1}}
\newcommand{\SpecialStringTok}[1]{\textcolor[rgb]{0.13,0.47,0.30}{#1}}
\newcommand{\StringTok}[1]{\textcolor[rgb]{0.13,0.47,0.30}{#1}}
\newcommand{\VariableTok}[1]{\textcolor[rgb]{0.07,0.07,0.07}{#1}}
\newcommand{\VerbatimStringTok}[1]{\textcolor[rgb]{0.13,0.47,0.30}{#1}}
\newcommand{\WarningTok}[1]{\textcolor[rgb]{0.37,0.37,0.37}{\textit{#1}}}

\providecommand{\tightlist}{%
  \setlength{\itemsep}{0pt}\setlength{\parskip}{0pt}}\usepackage{longtable,booktabs,array}
\usepackage{calc} % for calculating minipage widths
% Correct order of tables after \paragraph or \subparagraph
\usepackage{etoolbox}
\makeatletter
\patchcmd\longtable{\par}{\if@noskipsec\mbox{}\fi\par}{}{}
\makeatother
% Allow footnotes in longtable head/foot
\IfFileExists{footnotehyper.sty}{\usepackage{footnotehyper}}{\usepackage{footnote}}
\makesavenoteenv{longtable}
\usepackage{graphicx}
\makeatletter
\def\maxwidth{\ifdim\Gin@nat@width>\linewidth\linewidth\else\Gin@nat@width\fi}
\def\maxheight{\ifdim\Gin@nat@height>\textheight\textheight\else\Gin@nat@height\fi}
\makeatother
% Scale images if necessary, so that they will not overflow the page
% margins by default, and it is still possible to overwrite the defaults
% using explicit options in \includegraphics[width, height, ...]{}
\setkeys{Gin}{width=\maxwidth,height=\maxheight,keepaspectratio}
% Set default figure placement to htbp
\makeatletter
\def\fps@figure{htbp}
\makeatother

\usepackage{booktabs}
\usepackage{caption}
\usepackage{longtable}
\usepackage{colortbl}
\usepackage{array}
\usepackage{anyfontsize}
\usepackage{multirow}
\KOMAoption{captions}{tableheading}
\makeatletter
\@ifpackageloaded{caption}{}{\usepackage{caption}}
\AtBeginDocument{%
\ifdefined\contentsname
  \renewcommand*\contentsname{Table of contents}
\else
  \newcommand\contentsname{Table of contents}
\fi
\ifdefined\listfigurename
  \renewcommand*\listfigurename{List of Figures}
\else
  \newcommand\listfigurename{List of Figures}
\fi
\ifdefined\listtablename
  \renewcommand*\listtablename{List of Tables}
\else
  \newcommand\listtablename{List of Tables}
\fi
\ifdefined\figurename
  \renewcommand*\figurename{Figure}
\else
  \newcommand\figurename{Figure}
\fi
\ifdefined\tablename
  \renewcommand*\tablename{Table}
\else
  \newcommand\tablename{Table}
\fi
}
\@ifpackageloaded{float}{}{\usepackage{float}}
\floatstyle{ruled}
\@ifundefined{c@chapter}{\newfloat{codelisting}{h}{lop}}{\newfloat{codelisting}{h}{lop}[chapter]}
\floatname{codelisting}{Listing}
\newcommand*\listoflistings{\listof{codelisting}{List of Listings}}
\makeatother
\makeatletter
\makeatother
\makeatletter
\@ifpackageloaded{caption}{}{\usepackage{caption}}
\@ifpackageloaded{subcaption}{}{\usepackage{subcaption}}
\makeatother
\ifLuaTeX
  \usepackage{selnolig}  % disable illegal ligatures
\fi
\usepackage{bookmark}

\IfFileExists{xurl.sty}{\usepackage{xurl}}{} % add URL line breaks if available
\urlstyle{same} % disable monospaced font for URLs
\hypersetup{
  pdftitle={Capítulo I},
  pdfauthor={Ing. Marcelo Chávez},
  colorlinks=true,
  linkcolor={blue},
  filecolor={Maroon},
  citecolor={Blue},
  urlcolor={Blue},
  pdfcreator={LaTeX via pandoc}}

\title{Capítulo I}
\author{Ing. Marcelo Chávez}
\date{}

\begin{document}
\maketitle

\renewcommand*\contentsname{\textbf{CONTENIDOS}}
{
\hypersetup{linkcolor=green}
\setcounter{tocdepth}{2}
\tableofcontents
}
\chapter{Introducción}\label{introducciuxf3n}

\section{Las bases de datos de la
ENDI}\label{las-bases-de-datos-de-la-endi}

\subsection{Formulario de personas}\label{formulario-de-personas}

\begin{Shaded}
\begin{Highlighting}[]
\CommentTok{\# Conectar a PostgreSQL}
\NormalTok{postgresql\_conex }\OtherTok{\textless{}{-}} \FunctionTok{dbConnect}\NormalTok{(}\FunctionTok{PostgreSQL}\NormalTok{(),}
                 \AttributeTok{dbname =} \StringTok{"db\_stat"}\NormalTok{,}
                 \AttributeTok{host =} \StringTok{"localhost"}\NormalTok{, }
                 \AttributeTok{port =} \DecValTok{5432}\NormalTok{,}
                 \AttributeTok{user =} \StringTok{"postgres"}\NormalTok{,}
                 \AttributeTok{password =} \StringTok{"marce"}\NormalTok{)}

\CommentTok{\# Leer datos desde PostgreSQL}
\NormalTok{f1\_personas }\OtherTok{\textless{}{-}} \FunctionTok{dbGetQuery}\NormalTok{(postgresql\_conex, }\StringTok{"SELECT * FROM endi.f1\_personas"}\NormalTok{)}

\CommentTok{\# Lectura de la tabla de población menor de 5 años {-} censo 2022 INEC:}
\NormalTok{pob\_menor5anios }\OtherTok{\textless{}{-}} \FunctionTok{dbGetQuery}\NormalTok{(postgresql\_conex, }\StringTok{"SELECT * FROM endi.censo\_pob2022"}\NormalTok{)}

\CommentTok{\# Crear el objeto de diseño de encuesta}
\NormalTok{f1\_personas\_diseño }\OtherTok{\textless{}{-}}\NormalTok{ f1\_personas }\SpecialCharTok{\%\textgreater{}\%}
    \FunctionTok{as\_survey\_design}\NormalTok{(}\AttributeTok{ids =} \DecValTok{1}\NormalTok{, }\AttributeTok{weights =}\NormalTok{ fexp)}

\CommentTok{\# Calcular la prevalencia de desnutrición crónica por provincia}
\NormalTok{reporte\_dci }\OtherTok{\textless{}{-}}\NormalTok{ f1\_personas\_diseño }\SpecialCharTok{\%\textgreater{}\%}
    \FunctionTok{filter}\NormalTok{(f1\_s1\_3\_1 }\SpecialCharTok{\textless{}} \DecValTok{5}\NormalTok{) }\SpecialCharTok{\%\textgreater{}\%}
    \FunctionTok{group\_by}\NormalTok{(prov,provincia) }\SpecialCharTok{\%\textgreater{}\%}
    \FunctionTok{summarise}\NormalTok{(}
        \AttributeTok{indicador\_dci =} \FunctionTok{survey\_mean}\NormalTok{(dcronica }\SpecialCharTok{==} \DecValTok{1}\NormalTok{, }\AttributeTok{na.rm =} \ConstantTok{TRUE}\NormalTok{),  }\CommentTok{\# Indicador de DCI (\%)}
        \AttributeTok{total\_con\_dci =} \FunctionTok{round}\NormalTok{(}\FunctionTok{survey\_total}\NormalTok{(dcronica }\SpecialCharTok{==} \DecValTok{1}\NormalTok{, }\AttributeTok{na.rm =} \ConstantTok{TRUE}\NormalTok{), }\DecValTok{2}\NormalTok{),  }\CommentTok{\# Total niños con DCI}
        \AttributeTok{total\_ninios =} \FunctionTok{round}\NormalTok{(}\FunctionTok{survey\_total}\NormalTok{(fexp, }\AttributeTok{na.rm =} \ConstantTok{TRUE}\NormalTok{), }\DecValTok{2}\NormalTok{),  }\CommentTok{\# Total niños menores de 5 años}
        \AttributeTok{total\_sin\_dci =}\NormalTok{ total\_ninios }\SpecialCharTok{{-}}\NormalTok{ total\_con\_dci,  }\CommentTok{\# Total niños sin DCI}
        \AttributeTok{.groups =} \StringTok{"drop"}\NormalTok{) }\SpecialCharTok{\%\textgreater{}\%} 
    \FunctionTok{arrange}\NormalTok{(}\FunctionTok{desc}\NormalTok{(indicador\_dci)) }\SpecialCharTok{\%\textgreater{}\%} 
    \FunctionTok{inner\_join}\NormalTok{(pob\_menor5anios, }\AttributeTok{by=}\StringTok{"prov"}\NormalTok{) }\SpecialCharTok{\%\textgreater{}\%} 
    \FunctionTok{select}\NormalTok{(}\SpecialCharTok{{-}}\FunctionTok{ends\_with}\NormalTok{(}\StringTok{"\_se"}\NormalTok{), }\SpecialCharTok{{-}}\NormalTok{prov)}

\CommentTok{\# Generar la tabla interactiva}
\NormalTok{reporte\_dci }\SpecialCharTok{\%\textgreater{}\%}
    \FunctionTok{gt}\NormalTok{() }\SpecialCharTok{\%\textgreater{}\%}
    \FunctionTok{tab\_header}\NormalTok{(}
        \AttributeTok{title =} \FunctionTok{md}\NormalTok{(}\StringTok{"**Reporte de Desnutrición Crónica Infantil por Provincia**"}\NormalTok{),}
        \AttributeTok{subtitle =} \FunctionTok{md}\NormalTok{(}\StringTok{"Total de población y DCI en niños menores de 5 años"}\NormalTok{)}
\NormalTok{    ) }\SpecialCharTok{\%\textgreater{}\%}
    \FunctionTok{cols\_label}\NormalTok{(}
        \AttributeTok{provincia =} \StringTok{"Provincia"}\NormalTok{,}
        \AttributeTok{total\_ninios =} \StringTok{"Niños menores de 5 años {-} ENDI"}\NormalTok{,}
        \AttributeTok{total\_con\_dci =} \StringTok{"Niños con DCI {-} ENDI"}\NormalTok{,}
        \AttributeTok{total\_sin\_dci =} \StringTok{"Niños sin DCI {-} ENDI"}\NormalTok{,}
        \AttributeTok{indicador\_dci =} \StringTok{"Indicador DCI (\%)"}\NormalTok{,}
        \AttributeTok{pob\_menor5anios =} \StringTok{"Niños menores de 5 años {-} CPV2022"}
\NormalTok{    ) }\SpecialCharTok{\%\textgreater{}\%}
    \FunctionTok{fmt\_number}\NormalTok{(}
        \AttributeTok{columns =} \FunctionTok{vars}\NormalTok{(total\_ninios, total\_con\_dci, total\_sin\_dci, pob\_menor5anios),}
        \AttributeTok{decimals =} \DecValTok{0}\NormalTok{,}
        \AttributeTok{sep\_mark =} \StringTok{"."}
\NormalTok{    ) }\SpecialCharTok{\%\textgreater{}\%}
    \FunctionTok{fmt\_percent}\NormalTok{(}
        \AttributeTok{columns =}\NormalTok{ indicador\_dci,}
        \AttributeTok{decimals =} \DecValTok{1}
\NormalTok{    ) }\SpecialCharTok{\%\textgreater{}\%}
    \FunctionTok{tab\_options}\NormalTok{(}
        \AttributeTok{table.font.size =} \DecValTok{14}\NormalTok{,}
        \AttributeTok{table.border.top.color =} \StringTok{"black"}\NormalTok{,}
        \AttributeTok{table.border.bottom.color =} \StringTok{"black"}\NormalTok{,}
        \AttributeTok{table.border.left.color =} \StringTok{"black"}\NormalTok{,}
        \AttributeTok{table.border.right.color =} \StringTok{"black"}
\NormalTok{    ) }\SpecialCharTok{\%\textgreater{}\%}
    \CommentTok{\# Añadir iconos de flechas según el valor de DCI}
    \FunctionTok{text\_transform}\NormalTok{(}
        \AttributeTok{locations =} \FunctionTok{cells\_body}\NormalTok{(}\AttributeTok{columns =} \FunctionTok{vars}\NormalTok{(indicador\_dci)),}
        \AttributeTok{fn =} \ControlFlowTok{function}\NormalTok{(x) \{}
\NormalTok{            dplyr}\SpecialCharTok{::}\FunctionTok{case\_when}\NormalTok{(}
\NormalTok{                x }\SpecialCharTok{\textgreater{}=} \DecValTok{23} \SpecialCharTok{\textasciitilde{}} \FunctionTok{paste0}\NormalTok{(}\StringTok{"↑ "}\NormalTok{, x),  }\CommentTok{\# Nivel alto (flecha arriba)}
\NormalTok{                x }\SpecialCharTok{\textless{}} \DecValTok{23} \SpecialCharTok{\&}\NormalTok{ x }\SpecialCharTok{\textgreater{}} \DecValTok{16} \SpecialCharTok{\textasciitilde{}} \FunctionTok{paste0}\NormalTok{(}\StringTok{"⇄ "}\NormalTok{, x),  }\CommentTok{\# Nivel medio (flecha derecha)}
                \ConstantTok{TRUE} \SpecialCharTok{\textasciitilde{}} \FunctionTok{paste0}\NormalTok{(}\StringTok{"↓ "}\NormalTok{, x)  }\CommentTok{\# Nivel bajo (flecha abajo)}
\NormalTok{            )}
\NormalTok{        \}}
\NormalTok{    ) }\SpecialCharTok{\%\textgreater{}\%}
    \CommentTok{\# Mejorar la apariencia de las celdas de la tabla}
    \FunctionTok{tab\_style}\NormalTok{(}
        \AttributeTok{style =} \FunctionTok{list}\NormalTok{(}
            \FunctionTok{cell\_fill}\NormalTok{(}\AttributeTok{color =} \StringTok{"\#f5f5f5"}\NormalTok{),}
            \FunctionTok{cell\_text}\NormalTok{(}\AttributeTok{weight =} \StringTok{"bold"}\NormalTok{, }\AttributeTok{color =} \StringTok{"black"}\NormalTok{)}
\NormalTok{        ),}
        \AttributeTok{locations =} \FunctionTok{cells\_column\_labels}\NormalTok{()}
\NormalTok{    ) }\SpecialCharTok{\%\textgreater{}\%}
    \FunctionTok{tab\_style}\NormalTok{(}
        \AttributeTok{style =} \FunctionTok{list}\NormalTok{(}
            \FunctionTok{cell\_fill}\NormalTok{(}\AttributeTok{color =} \StringTok{"\#e0f7fa"}\NormalTok{)}
\NormalTok{        ),}
        \AttributeTok{locations =} \FunctionTok{cells\_body}\NormalTok{(}\AttributeTok{columns =} \FunctionTok{vars}\NormalTok{(provincia))}
\NormalTok{    )}
\end{Highlighting}
\end{Shaded}

\begin{verbatim}
Warning: Since gt v0.3.0, `columns = vars(...)` has been deprecated.
* Please use `columns = c(...)` instead.
Since gt v0.3.0, `columns = vars(...)` has been deprecated.
* Please use `columns = c(...)` instead.
Since gt v0.3.0, `columns = vars(...)` has been deprecated.
* Please use `columns = c(...)` instead.
\end{verbatim}

\begin{verbatim}
Warning in prettyNum(r, big.mark = big.mark, big.interval = big.interval, :
'big.mark' and 'decimal.mark' are both '.', which could be confusing
Warning in prettyNum(r, big.mark = big.mark, big.interval = big.interval, :
'big.mark' and 'decimal.mark' are both '.', which could be confusing
Warning in prettyNum(r, big.mark = big.mark, big.interval = big.interval, :
'big.mark' and 'decimal.mark' are both '.', which could be confusing
Warning in prettyNum(r, big.mark = big.mark, big.interval = big.interval, :
'big.mark' and 'decimal.mark' are both '.', which could be confusing
\end{verbatim}

\begin{table}
\caption*{
{\large \textbf{Reporte de Desnutrición Crónica Infantil por Provincia}} \\ 
{\small Total de población y DCI en niños menores de 5 años}
} 
\fontsize{10.5pt}{12.6pt}\selectfont
\begin{tabular*}{\linewidth}{@{\extracolsep{\fill}}lrrrrr}
\toprule
Provincia & Indicador DCI (\%) & Niños con DCI - ENDI & Niños menores de 5 años - ENDI & Niños sin DCI - ENDI & Niños menores de 5 años - CPV2022 \\ 
\midrule\addlinespace[2.5pt]
{\cellcolor[HTML]{E0F7FA}{CHIMBORAZO}} & ↑ 34.1\% & 188 & 635 & 447 & 32.462 \\ 
{\cellcolor[HTML]{E0F7FA}{SANTA ELENA}} & ↑ 30.0\% & 200 & 737 & 537 & 34.839 \\ 
{\cellcolor[HTML]{E0F7FA}{COTOPAXI}} & ↑ 29.4\% & 193 & 1.221 & 1.028 & 35.693 \\ 
{\cellcolor[HTML]{E0F7FA}{PASTAZA}} & ↑ 27.4\% & 57 & 104 & 48 & 12.158 \\ 
{\cellcolor[HTML]{E0F7FA}{TUNGURAHUA}} & ↑ 27.0\% & 184 & 445 & 261 & 38.016 \\ 
{\cellcolor[HTML]{E0F7FA}{MORONA SANTIAGO}} & ↑ 27.0\% & 97 & 333 & 236 & 23.095 \\ 
{\cellcolor[HTML]{E0F7FA}{BOLÍVAR}} & ↑ 25.4\% & 57 & 73 & 17 & 14.717 \\ 
{\cellcolor[HTML]{E0F7FA}{CAÑAR}} & ↑ 23.0\% & 59 & 96 & 37 & 18.013 \\ 
{\cellcolor[HTML]{E0F7FA}{CARCHI}} & ⇄ 22.9\% & 54 & 103 & 49 & 11.341 \\ 
{\cellcolor[HTML]{E0F7FA}{IMBABURA}} & ⇄ 21.7\% & 143 & 831 & 688 & 31.123 \\ 
{\cellcolor[HTML]{E0F7FA}{PICHINCHA}} & ⇄ 20.6\% & 684 & 7.000 & 6.316 & 184.640 \\ 
{\cellcolor[HTML]{E0F7FA}{AZUAY}} & ⇄ 20.2\% & 168 & 866 & 698 & 54.110 \\ 
{\cellcolor[HTML]{E0F7FA}{ORELLANA}} & ⇄ 17.6\% & 59 & 249 & 190 & 19.273 \\ 
{\cellcolor[HTML]{E0F7FA}{NAPO}} & ⇄ 17.3\% & 38 & 76 & 38 & 13.080 \\ 
{\cellcolor[HTML]{E0F7FA}{LOJA}} & ⇄ 16.6\% & 102 & 814 & 712 & 34.367 \\ 
{\cellcolor[HTML]{E0F7FA}{MANABÍ}} & ⇄ 16.2\% & 372 & 6.930 & 6.558 & 130.150 \\ 
{\cellcolor[HTML]{E0F7FA}{ZAMORA CHINCHIPE}} & ↓ 14.5\% & 24 & 43 & 18 & 10.418 \\ 
{\cellcolor[HTML]{E0F7FA}{GUAYAS}} & ↓ 12.7\% & 695 & 36.680 & 35.985 & 342.644 \\ 
{\cellcolor[HTML]{E0F7FA}{ESMERALDAS}} & ↓ 12.4\% & 149 & 1.759 & 1.609 & 55.478 \\ 
{\cellcolor[HTML]{E0F7FA}{SUCUMBÍOS}} & ↓ 11.7\% & 43 & 213 & 170 & 17.740 \\ 
{\cellcolor[HTML]{E0F7FA}{LOS RÍOS}} & ↓ 11.5\% & 170 & 3.637 & 3.467 & 81.018 \\ 
{\cellcolor[HTML]{E0F7FA}{SANTO DOMINGO DE LOS TSÁCHILAS}} & ↓ 11.1\% & 68 & 396 & 329 & 41.032 \\ 
{\cellcolor[HTML]{E0F7FA}{EL ORO}} & ↓ 10.1\% & 82 & 559 & 477 & 55.901 \\ 
\bottomrule
\end{tabular*}
\end{table}

\begin{Shaded}
\begin{Highlighting}[]
\CommentTok{\# generar\_tabla\_dt(}
\CommentTok{\#   data = prevalencia\_dci\_prov,}
\CommentTok{\#   colnames = c("Provincia", "DCI (\%)", "Prevalencia SE", "Total DCI", "Total DCI SE", "Total Niños", "Total Niños SE"),}
\CommentTok{\#   title = "📊 Prevalencia de Desnutrición Crónica Infantil a Nivel Provincial",}
\CommentTok{\#   subtitle = "Fuente: Encuesta Nacional de Desnutrición Infantil (ENDI) {-} INEC, 2023.",}
\CommentTok{\#   footnote = "Datos oficiales del INEC, procesamiento 2024.",}
\CommentTok{\#   percentage\_col = "prevalencia",}
\CommentTok{\#   header\_color = "\#004488",}
\CommentTok{\#   text\_color = "\#FFFFFF",}
\CommentTok{\#   border\_color = "\#CCCCCC",}
\CommentTok{\#   height = 700)}
\end{Highlighting}
\end{Shaded}

\chapter{Explicación por provincia sobre los porcentajes de DCI en niños
menores de 5 años (CPV
2022)}\label{explicaciuxf3n-por-provincia-sobre-los-porcentajes-de-dci-en-niuxf1os-menores-de-5-auxf1os-cpv-2022}

\begin{itemize}
\tightlist
\item
  \textbf{Chimborazo (34.1\%)}: Presenta el porcentaje más alto de
  desnutrición crónica infantil, reflejando una grave problemática de
  acceso a una adecuada alimentación y atención nutricional.\\
\item
  \textbf{Santa Elena (30.0\%)}: Alto nivel de DCI, posiblemente
  asociado a limitaciones en acceso a servicios de salud y alimentación
  balanceada.\\
\item
  \textbf{Cotopaxi (29.4\%)}: Similar a Chimborazo, su alta prevalencia
  de DCI puede estar relacionada con condiciones socioeconómicas y
  acceso a recursos básicos.\\
\item
  \textbf{Pastaza (27.4\%)}: La desnutrición en esta provincia amazónica
  podría estar influenciada por barreras geográficas y acceso limitado a
  servicios de salud.\\
\item
  \textbf{Tungurahua (27.0\%)}: A pesar de ser una provincia con
  desarrollo económico, la DCI sigue siendo un problema, especialmente
  en zonas rurales.\\
\item
  \textbf{Morona Santiago (27.0\%)}: La desnutrición infantil en
  comunidades amazónicas sigue siendo un desafío por factores
  socioeconómicos y de acceso.\\
\item
  \textbf{Bolívar (25.4\%)}: Provincia con altos niveles de DCI,
  reflejando desafíos en nutrición y atención infantil.\\
\item
  \textbf{Cañar (23.0\%)}: Aunque con un porcentaje menor que otras
  provincias de la Sierra, la DCI sigue siendo preocupante.\\
\item
  \textbf{Carchi (22.9\%)}: La desnutrición en esta provincia puede
  estar influenciada por factores económicos y disponibilidad de
  alimentos.\\
\item
  \textbf{Imbabura (21.7\%)}: Se mantiene dentro de los niveles medios
  de DCI, aunque sigue siendo un problema relevante.\\
\item
  \textbf{Pichincha (20.6\%)}: A pesar de ser la provincia con mayor
  desarrollo económico, aún hay sectores vulnerables con DCI.\\
\item
  \textbf{Azuay (20.2\%)}: La DCI en esta provincia puede estar
  relacionada con desigualdades en el acceso a servicios.\\
\item
  \textbf{Orellana (17.6\%)}: Menor porcentaje en comparación con otras
  provincias amazónicas, pero sigue siendo un problema de salud
  pública.\\
\item
  \textbf{Napo (17.3\%)}: Similar a Orellana, con una tasa considerable
  de DCI en niños menores de 5 años.\\
\item
  \textbf{Loja (16.6\%)}: Provincia con menor DCI en la Sierra, aunque
  la problemática persiste en sectores rurales.\\
\item
  \textbf{Manabí (16.2\%)}: En la Costa, esta provincia presenta uno de
  los niveles más bajos de DCI.\\
\item
  \textbf{Zamora Chinchipe (14.5\%)}: La provincia con menor DCI en la
  Amazonía, posiblemente por mejores condiciones de acceso a
  alimentos.\\
\item
  \textbf{Guayas (12.7\%)}: Aunque tiene una baja prevalencia en
  comparación con la Sierra, aún existen sectores vulnerables con
  desnutrición.\\
\item
  \textbf{Esmeraldas (12.4\%)}: Similar a Guayas, con una tasa menor de
  DCI en la Costa.\\
\item
  \textbf{Sucumbíos (11.7\%)}: Provincia amazónica con menor prevalencia
  de DCI en comparación con sus vecinas.\\
\item
  \textbf{Los Ríos (11.5\%)}: Presenta uno de los niveles más bajos de
  DCI en la Costa.\\
\item
  \textbf{Santo Domingo de los Tsáchilas (11.1\%)}: Provincia con una de
  las tasas más bajas de DCI.\\
\item
  \textbf{El Oro (10.1\%)}: La provincia con el menor porcentaje de DCI
  en el país, posiblemente debido a una mejor alimentación y acceso a
  servicios de salud.
\end{itemize}

Este análisis revela que las provincias de la \textbf{Sierra Central y
la Amazonía} presentan los niveles más altos de desnutrición crónica
infantil (DCI), lo que indica que estos territorios enfrentan desafíos
estructurales en términos de acceso a servicios de salud, seguridad
alimentaria y condiciones socioeconómicas. La presencia de comunidades
indígenas en estas zonas también juega un papel clave, ya que
históricamente han enfrentado desigualdades en la provisión de servicios
básicos y políticas de desarrollo específicas. En particular, provincias
como \textbf{Chimborazo, Cotopaxi y Bolívar} han registrado
históricamente cifras alarmantes de DCI, reflejando la persistencia del
problema a lo largo del tiempo.

A nivel \textbf{nacional}, la desnutrición crónica infantil sigue siendo
un problema prioritario en Ecuador. A pesar de los esfuerzos
gubernamentales y de organizaciones internacionales para reducir estos
índices, los datos reflejan que aún persisten brechas significativas en
las políticas de alimentación y salud infantil. Factores como la
pobreza, el acceso limitado a agua potable y saneamiento, y la falta de
educación nutricional en los hogares continúan incidiendo en los altos
niveles de DCI. Además, los efectos de la crisis económica y la pandemia
de COVID-19 han exacerbado las dificultades para mejorar los indicadores
de nutrición infantil en el país.

Desde una perspectiva \textbf{regional}, Ecuador no es un caso aislado.
Países vecinos como \textbf{Perú y Bolivia} también han registrado
históricamente altos niveles de DCI, especialmente en sus regiones
andinas y rurales. Sin embargo, experiencias exitosas en estos países
han demostrado que programas de intervención focalizada, como
transferencias condicionadas de dinero, campañas de educación
nutricional y la promoción de la lactancia materna exclusiva, pueden
reducir significativamente la DCI. Ecuador podría beneficiarse de estas
experiencias mediante la implementación de políticas más integrales y
sostenibles.

En contraste, las provincias de la \textbf{Costa}, como \textbf{Guayas,
Manabí y Esmeraldas}, presentan niveles más bajos de DCI en comparación
con la Sierra y la Amazonía. Esto podría explicarse por una mayor
disponibilidad de alimentos debido a la producción agrícola y pesquera,
así como un acceso relativamente mejor a los mercados y centros urbanos
con infraestructura hospitalaria más desarrollada. No obstante, en estas
provincias todavía existen sectores vulnerables, especialmente en
comunidades rurales y barrios marginales de las grandes ciudades, donde
la pobreza y la inseguridad alimentaria continúan siendo problemáticas
latentes.

Para abordar este desafío, es crucial que Ecuador fortalezca sus
estrategias de intervención mediante un enfoque intersectorial que
involucre a ministerios, gobiernos locales, sociedad civil y organismos
internacionales. La implementación de programas que integren el acceso a
una alimentación adecuada, mejoras en los servicios de salud y campañas
de concienciación pueden marcar una diferencia significativa en la
reducción de la DCI. Asimismo, el monitoreo constante de los indicadores
y la inversión en nutrición infantil deben ser prioridades dentro de la
agenda de desarrollo del país.



\end{document}
